
%%%%%%%%%%%%%%%%%%%%%%%%%%%%%%%%%%%%%%%%
%           main body
%%%%%%%%%%%%%%%%%%%%%%%%%%%%%%%%%%%%%%%%

%-----------------------------------
%	TITLE PAGE
%-----------------------------------

\title[数学分析]{\texorpdfstring{\LARGE \bfseries 第九章\quad 数项级数}{第九章 数项级数}} % The short title appears at the bottom of every slide, the full title is only on the title page
\author[正项级数] % Your institution as it will appear on the bottom of every slide, maybe shorthand to save space
{\Large \bfseries 正项级数}
% \author{Singular Value Decomposition} % Your name
% \date{\today} % Date, can be changed to a custom date
\date{}

%------------------------------------------------

\begin{document}

%------------------------------------------------

\begin{frame}

\titlepage % Print the title page as the first slide

\end{frame}

%------------------------------------------------

% \begin{frame}
% \frametitle{内容提要} % Table of contents slide, comment this block out to remove it
% \tableofcontents % Throughout your presentation, if you choose to use \section{} and \subsection{} commands, these will automatically be printed on this slide as an overview of your presentation
% \end{frame}

%-----------------------------------
%	PRESENTATION SLIDES
%-----------------------------------

%------------------------------------------------

\begin{frame}{正项级数比较判别法存在的问题与改进方案}

用比较判别法时，先要对所考虑的级数的收敛性有一个大致估计，进而通过放缩，寻找一个{\color{red}敛散性已知或相对容易判断}的合适级数与之项比较。但就绝大多数情况而言，这两个步骤都具有相当的难度，因此，我们希望有一些
\begin{empheq}[box={\fancymath[colback=cyan!30,drop lifted shadow, sharp corners]}]{equation*}
  \text{\bfseries\color{red!80} 仅针对级数自身元素 (通项) 进行分析的敛散性判别方法。}
\end{empheq}

\pause

正项等比级数$\sum\limits_{n=1}^\infty q^n ~ (q > 0)$给我们很重要的启示：它的敛散性只依赖于其相邻两项之比，即公比$q$是否小于$1.$ 一个自然的想法是，如果

\end{frame}

%------------------------------------------------

\section[柯西判别法]{正项级数敛散性的柯西判别法}

%------------------------------------------------

\begin{frame}{待写}



\end{frame}

%------------------------------------------------

\section[达朗贝尔判别法]{正项级数敛散性的达朗贝尔判别法}

%------------------------------------------------

\begin{frame}{待写}



\end{frame}